\icode[magenta]{N21-MD-MU-L4-261112-001}
\nceeID[6]{2021全国乙卷}{21}
水平地面上有一质量为\(m_1\)的长木板，木板的左端上有一质量为\(m_2\)的物块，如图\((a)\)所示。用水平向右的拉力\(F\)作用在物块上，\(F\)随时间\(t\)的变化关系如图\((b)\)所示，其中\(F_1\)、\(F_2\)分别为\(t_1\)、\(t_2\)时刻\(F\)的大小。木板的加速度\(a_1\)随时间\(t\)的变化关系如图\((c)\)所示。已知木板与地面间的动摩擦因数为\(\mu_1\)，物块与木板间的动摩擦因数为\(\mu_2\)。假设最大静摩擦力均与相应的滑动摩擦力相等，重力加速度大小为\(g\)。则\paren[]

\begin{tikzpicture}[>=stealth,font=\scriptsize]
	\draw[fill=teal!60] (0,0) rectangle(3,.5)node[midway]{\(m_1\)};
	\draw[fill=orange!60] (0,.5) rectangle(.6,.9)node[midway]{\(m_2\)};
	\foreach \x in {-1,-.8,...,3.6} 
		\draw[gray] (\x,-.2)--++(.2,.2);
	\draw[thick] (-1,0)--(3.8,0);
	\draw[->] (.6,.7)--(1.6,.7)node[right]{\(F\)};
	\node at (1.3,-.8){图（\(a\)）};
\end{tikzpicture}
\begin{tikzpicture}[>=stealth,font=\scriptsize]
	\begin{scope}
		\draw[->,blue] (0,0)node[below left]{\(O\)}--(3,0)node[below]{\(t\)};
		\draw[->,blue] (0,0)--(0,2.5)node[right]{\(F\)};
		\draw[magenta] (0,0)--(2.5,2);
		\draw[dashed,cyan] (1,0)node[below,blue]{\(t_1\)}--(1,.8)--(0,.8)node[left,blue]{\(F_1\)} (1.5,0)node[below,blue]{\(t_2\)}--(1.5,1.2)--(0,1.2)node[left,blue]{\(F_2\)};
		\node at (1.5,-.8){图（\(b\)）};
	\end{scope}
	\begin{scope}[xshift=4cm]
		\draw[->,blue] (0,0)node[below left]{\(O\)}--(3,0)node[below]{\(t\)};
		\draw[->,blue] (0,0)--(0,2.5)node[right]{\(a_1\)};
		\draw[magenta] (1,0)node[below,blue]{\(t_1\)}--(1.5,1.6)--(2.2,1.6);
		\draw[dashed,cyan] (1.5,0)node[below,blue]{\(t_2\)}--(1.5,1.6)--(0,1.6);
		\node at (1.5,-.8){图（\(c\)）};
	\end{scope}	
\end{tikzpicture}	
\begin{choices}
	\item \(F_1 = \mu_1m_1g\)
	\item \(F_2 = \dfrac{m_2(m_1 + m_2)}{m_1}(\mu_2 -\mu_1)g\)
	\item \(\mu_2 > \dfrac{m_1 + m_2}{m_2}\mu_1\)
	\item 在\(0\sim t_2\)时间段物块与木板加速度相等
\end{choices}
\begin{answer}
    
\end{answer}